\chapter{Group actions, characters and block parameters}
\label{chap:library-current}

Equivariant maps and character identities often determine an action without
requiring an explicit choice of representatives. We study this principle for
permutation modules, block stabilisers and ordinary characters over different
fields. We also describe the quotient associated with a similitude multiplier
and the relation between block idempotents and rational series.

\section{Fixed bases and permutation modules}

Let $A$ act linearly on a module $V$ over a nonzero ring $R$. If $A$ fixes every
vector of an $R$-basis, it acts trivially on $V$: a linear map is determined
by its values on a basis. If another basis is indexed by an $A$-set $Y$
and the action permutes its basis vectors in accordance with the action on
$Y$, the linear independence of that basis implies that $A$ fixes $Y$
pointwise.

\begin{proposition}
Let $X$ and $Y$ be finite $A$-sets. Suppose that
$d\colon\Z X\longrightarrow\Z Y$ is an $A$-equivariant isomorphism of
free abelian groups. If $A$ fixes $X$ pointwise, it fixes $Y$ pointwise,
and there is an $A$-equivariant bijection $X\longrightarrow Y$.
\end{proposition}

\begin{proof}
The images under $d$ of the basis vectors indexed by $X$ form a pointwise
fixed basis of $\Z Y$. The preceding argument shows that $A$ acts trivially
on $\Z Y$, and hence on $Y$. Equality of the ranks gives $|X|=|Y|$.
Any bijection between the two finite sets is then equivariant.
\end{proof}

Let $X$ be an $A$-stable integral basic set of a block $b$, and suppose
that the restriction of the decomposition homomorphism gives an
$A$-equivariant isomorphism $\Z X\longrightarrow\Z\IBr(b)$.
If $A$ fixes every character in $X$, the proposition shows that it
fixes every irreducible Brauer character of $b$ and gives an
$A$-equivariant bijection $X\longrightarrow\IBr(b)$.

\begin{implementation}
\module{ManuscriptIBAW.Library.FixedBasis} specialises Mathlib's basis
extensionality, injectivity and equivalence of basis index sets to the
stated actions. \module{ManuscriptIBAW.Library.FixedBasicSet} applies
this argument to Mathlib's integral permutation modules. The block application is
\module{ManuscriptIBAW.TypeC.EvenBasicSet}.
\end{implementation}

\section{Equivariant bijections for involutions}

\begin{proposition}
Let $A$ be a group acting on finite sets $X$ and $Y$. Suppose that $A$ is generated by a
set $T$ and an element $\delta$, that every element of $T$ fixes both sets
pointwise, and that $\delta$ induces an involution on each set. There is
an $A$-equivariant bijection $X\longrightarrow Y$ if
\[
 |X|=|Y|\qquad\text{and}\qquad |X^{\langle\delta\rangle}|
   =|Y^{\langle\delta\rangle}|.
\]
\end{proposition}

\begin{proof}
Each orbit under $\delta$ has size one or two. The two equalities give
equal numbers of orbits of each size. Match the fixed points and the
orbits of size two to obtain a bijection commuting with $\delta$.
It commutes with $T$ because $T$ acts trivially. The elements of $A$
with which the bijection commutes form a subgroup, so the generating
hypothesis gives equivariance under $A$.
\end{proof}

The hypothesis concerns the permutations induced by $\delta$. It does
not require that $\delta$ itself have order two in $A$. This distinction
allows the proposition to be used when $\delta^2$ is an inner automorphism
and the objects under consideration are conjugacy classes.

\begin{implementation}
\module{ManuscriptIBAW.Library.EquivariantBijections} proves the subgroup
and generating set arguments and uses
\lean{Formalisation.C2Cancellation} for the two involutions. That result
specialises Mathlib's conjugacy classification of finite permutations by cycle type
and its fixed point count. The local statements express the conclusion as
an equivariant bijection for the given actions.
\end{implementation}

\section{Independent signs on a direct sum}

Let $V=\bigoplus_{c\in I}V_c$ be a finite dimensional space over a field of
characteristic different from two. For each $c$, let $\varepsilon_c$ act
as $-1$ on $V_c$ and as $1$ on the other summands. A linear map commuting
with every $\varepsilon_c$ preserves every $V_c$. Indeed, in these
coordinates, commutation forces an entry between distinct summands to
equal its negative. That entry is zero because two is invertible.

If the decomposition is orthogonal for a nondegenerate alternating form
and the map is symplectic, its restriction to each summand is symplectic.
The form equation restricts to each diagonal block. Consequently the map
belongs to the product of the symplectic groups of the summands.

Suppose that $R=\prod_c R_c$ acts independently on these summands, that
$-1\in R_c$, and that the centraliser of $R_c$ in its symplectic group
is contained in $R_c$. Every element centralising $R$ commutes with the
independent signs, hence preserves each summand. Its restriction there
centralises $R_c$ and therefore belongs to $R_c$. Thus $C_{\Sp(V)}(R)\leq R$.
The index $c$ distinguishes different summands, including copies on which
the same group acts.

\begin{implementation}
\lean{ManuscriptIBAW.TypeC.independentSignSeparation} in
\module{ManuscriptIBAW.TypeC.PrincipalProducts} proves the matrix argument
for the specified symplectic coordinates over a finite field of odd order.
Mathlib supplies matrix reindexing and the block diagonal multiplication
and transpose identities. The proof combines them with the specified
embedding and Gram matrix equations to separate the summands.
\lean{ManuscriptIBAW.TypeC.PrincipalProductData.signs} is a proved consequence
of those equations. The local centraliser assertion remains a source
assumption in the application.
\end{implementation}

\section{Stabilisers under an equivariant bijection}

Let $A=M\rtimes E$ act on sets $X$ and $Y$, and let
$\beta:X\longrightarrow Y$ be an $A$-equivariant bijection. Then
$A_x=A_{\beta(x)}$, since equivariance and injectivity give
$\beta(ax)=\beta(x)$ if and only if $ax=x$. In particular, if
\[
 m(e\beta(x))=\beta(x)
 \quad\Longleftrightarrow\quad
 m\beta(x)=\beta(x)\ \text{and}\ e\beta(x)=\beta(x),
\]
the same equivalence holds with $x$ in place of $\beta(x)$. This transfers
the factorisation of a stabiliser into its $M$ and $E$ parts.

Chapter~\ref{chap:library:assembly} constructs a global bijection from
bijections over representatives of block orbits that are equivariant under
their stabilisers.
It preserves the block assignments. For character sets of a finite group
$G$, suppose that $A$ acts through the usual automorphism action and maps
onto $\Aut(G)$. Equivariance under $A$ then gives equivariance under every
automorphism: choose a preimage and apply the same identity.

\begin{implementation}
\module{ManuscriptIBAW.TypeC.OddGlobalFactorization} constructs
\lean{globalBijection} from the block basic sets, the stated quotient
actions and the Conlon and Burnside hypotheses. In that namespace,
\lean{GlobalBijection.automorphism_equivariant} proves full automorphism
equivariance, and \lean{GlobalBijection.allBrauerFactorization} transfers
ordinary stabiliser factorisation through that same bijection. The
construction uses the basic set hypotheses first, then transfers the
factorisation to the Brauer characters.
\end{implementation}

\section{Linear characters and block stabilisers}

Under the coefficient and tensor compatibility hypotheses of
Section~\ref{lib:blocks:linear-stabilisers}, let $G$ be finite,
let $\ell$ be prime and let $N\lhd G$.
Let $b$ be an $\ell$-block of $G$, and let $C$ be a
subgroup of the group of ordinary linear characters of $G/N$ of order
prime to $\ell$. Inflate these characters to $G$ and let them act by
tensoring. Each character in $C_b$ is trivial on $Z(G)$, so it factors
through $G/NZ(G)$. In particular,
\[
 |C_b|\leq |G:NZ(G)|.
\]
The same assertion holds for the corresponding linear Brauer characters.
The proof uses the common central character of the irreducible Brauer
characters of $b$. A character stabilising $b$ is trivial on the
$\ell'$-part of $Z(G)$ by this central character, and on its $\ell$-part
because its order is prime to $\ell$.

Now let a finite cyclic group $E$ act on $C$, and suppose that $C\rtimes E$
acts on the blocks in a way whose restriction to $C$ is the tensor
action just described. For $J=(C\rtimes E)_b$, let $D$ be the kernel
of the projection $J\longrightarrow E$. Sending an element of $D$ to
its $C$ coordinate gives an injection $D\longrightarrow C_b$.
Thus $|D|\leq |G:NZ(G)|$. If this index is at most two, $D$ is a normal
$2$-subgroup of $J$ and $J/D$ is cyclic.

Let $P/D$ be the Sylow $2$-subgroup of $J/D$. Then $P\lhd J$ is a
$2$-group and $J/P$ is cyclic of odd order. For every subgroup $U\leq J$,
the subgroup $U\cap P$ is a normal $2$-subgroup and
$U/(U\cap P)$ embeds in $J/P$. Consequently $U/O_2(U)$ is cyclic of
odd order. Thus every subgroup of $J$ is $2$-hypoelementary, as required
in the application of Conlon's theorem.

\begin{implementation}
The central character argument is in
\module{ModularRep.LinearBlockStabilizer}.
\module{ManuscriptIBAW.Library.TensorStabilizer} identifies the kernel of the
projection to $E$ with its character coordinate. It combines the resulting
bound with the hypoelementary group calculation in ModularRep.
Compatibility of the block action with tensoring
simple modules remains an explicit hypothesis.
\end{implementation}

\newpage
\begin{samepage}
\section{The multiplier quotient}

\begin{proposition}
Let $\mu\colon G\longrightarrow U$ be a surjective homomorphism to an
abelian group. It induces an isomorphism
\[
 G/(\ker\mu\,Z(G))\ \cong\ U/\mu(Z(G)).
\]
If $U=\F_q^\times$ and $\mu(Z(G))=(\F_q^\times)^2$, this quotient has
order at most two, and has order two when $q$ is odd.
\end{proposition}
\end{samepage}

\begin{proof}
The composite $G\longrightarrow U\longrightarrow U/\mu(Z(G))$ is
surjective with kernel $\ker\mu\,Z(G)$, so the first assertion is the
first isomorphism theorem. In $\F_q^\times$, the kernel of the squaring
map is the group of roots of $x^2=1$. It has at most two elements and
has exactly two in odd characteristic. The index of the image of an
endomorphism of a finite group equals the order of its kernel.
\end{proof}

For the conformal symplectic group of positive rank, the multiplier is specified by
$\langle gv,gw\rangle=\mu(g)\langle v,w\rangle$. Its kernel is the
symplectic subgroup. Surjectivity and the image of the centre are the
structural hypotheses in this application. The proposition determines the quotient index
from these hypotheses.

\begin{implementation}
\module{ManuscriptIBAW.Library.MultiplierQuotient} obtains the isomorphism
by identifying the kernel of the composite and applying Mathlib's first
isomorphism theorem. The finite field calculation combines Mathlib's
index formula for an endomorphism with its bounds and exact counts for
roots of unity.
\module{ManuscriptIBAW.TypeC.ConformalMultiplier} applies them to the
matrix groups with the specified symplectic form.
\end{implementation}

\section{Ordinary characters over a common cyclotomic field}

Let $G$ be finite, let $N$ be a positive multiple of $|G|$, and set
$L=\Q(\zeta_N)$. Fix embeddings $\iota_{\C}\colon L\longrightarrow\C$
and $\iota_K\colon L\longrightarrow K$, where $K$ has characteristic
zero. For character independence and the criterion for realisability,
see, for example, Serre, Section 12.1, Propositions 32--33. The cyclotomic
realisation theorem appears there as Theorem 24 (Brauer), with its
corollary, in Section 12.3 \cite{Serre1977}.

Assume that a realisation statement supplies, for each complex
irreducible character $\chi$, characters $\chi_L\in\Irr_L(G)$ and
$\chi_K\in\Irr_K(G)$. If the chosen basic set associates a simple $K[G]$-module $V_\chi$ with
$\chi$, the character of its representation $\rho_\chi$ must equal $\chi_K$. Thus the required identities are
\[
 \iota_{\C}(\chi_L(g))=\chi(g),\qquad
 \iota_K(\chi_L(g))=\chi_K(g)=\tr_K(\rho_\chi(g))
 \qquad(g\in G).
\]
These identities relate the two realisations to the selected labels.
They can also be obtained by constructing those labels from the
realisation statement.

Suppose that a group of automorphisms permutes a chosen set of complex
characters. Injectivity of $\iota_{\C}$ transports the character
identities to $L$, and applying $\iota_K$ gives the same identities
over $K$. Equality of the traces of the corresponding irreducible
representations gives their isomorphism over $K$, and hence equality
of their classes in the exact Grothendieck group. Thus the action on
the ordinary labels used by the decomposition homomorphism follows
from the action on the complex characters and the stated realisations.

The argument uses only the two embeddings of $L$, so it requires neither
an embedding of $\C$ into $K$ nor algebraic closure of $K$.
Existence of the realisations is an assumption. Their compatibility with
automorphisms follows from these trace identities.

\begin{implementation}
The common field is Mathlib's cyclotomic field. Trace separation uses
Mathlib's character averaging formula for the dimension of the intertwiner
space and its theorem that a nonzero intertwiner between irreducibles is
bijective. ModularRep combines these results over the stated field $K$.
In \module{ManuscriptIBAW.Characters.CyclotomicOrdinaryLabels},
\lean{OrdinaryLabelRealisation.ordinary_value} is the trace identity.
For labels obtained from source realisations,
\lean{SerreRealisationSource.labelRealisation} proves it.
\end{implementation}

\section{Primitive blocks and rational series}

Let $(b_B)_{B\in\mathcal B}$ be a finite complete family of nonzero
orthogonal central idempotents of a ring. Assign to each $B$ a parameter
$t_B$ in a finite group $H^*$. For $t\in H^*$, define
\[
 e_t=\sum_{\substack{B\in\mathcal B\\t_B\sim t}} b_B,
\]
where $\sim$ denotes conjugacy in $H^*$. Orthogonality gives
\[
 b_Be_t=
 \begin{cases}
 b_B,&t_B\sim t,\\
 0,&t_B\not\sim t.
 \end{cases}
\]
Since $b_B\ne0$, we have $b_Be_t=b_B$ if and only if $t_B$ is
conjugate to $t$.

In Jordan reduction the idempotents are the blocks of the
specified group algebra. Each $t_B$ is required to be semisimple and of
order prime to $\ell$. Let $s$ index a geometric series, and write $t_s\in H^*$ for its
assigned element of the finite dual group. The geometric series idempotent at $s$ is identified
with this same sum $e_{t_s}$. Completeness of the chosen semisimple
parameters then places every block in one of the series.

The label used to test strict quasi-isolation is also identified with
the image of $t_B$ under a specified injective homomorphism from $H^*$
to the algebraic dual group. The intended map includes rational points
in the algebraic group. Compatibility of series with automorphisms is a
separate identity for the induced map of group algebras and the dual action.
These identifications specify the blocks, labels and action to which
Jordan reduction applies. Their interpretation as rational series in
algebraic groups remains an assumption.

\begin{implementation}
\module{ManuscriptIBAW.Jordan.SeriesBinding} defines
\lean{rationalSeriesIdempotent} and the identification record
\lean{SeriesBinding} using the block family in ModularRep.
Mathlib supplies the group algebra and finite sum operations. The local
calculation applies orthogonality to this specified sum, while its
identification with the geometric series remains a hypothesis.
The multiplication formula is
\lean{block_mul_rationalSeriesIdempotent}. The membership equivalence and
coverage of the blocks are \lean{SeriesBinding.block_belongs_iff}
and \lean{SeriesBinding.every_block_belongs}, all in
\lean{ManuscriptIBAW.Jordan}.
\end{implementation}

